% Generated by scripts/assemble.py; corrections belong in source/corrections.json.
\examhead{185}{长安大学 2025 年数学分析试卷}{174}
% p175-b15; source PDF page 175
\par\Needspace{5\baselineskip}
\blockanchor{p175-b15}
\majorquestion{一}计算题. 每题 5 分, 共 30 分.\par
% p175-b16; source PDF page 175
\blockanchor{p175-b16}
\question{1}作废.\par
% p175-b17; source PDF page 175
\blockanchor{p175-b17}
\question{2}求数列极限 $\lim_{n\to \infty}n\left[\left(1 + \frac{1}{n}\right)^n -e\right].$\par
% p175-b18; source PDF page 175
\blockanchor{p175-b18}
\question{3}设 $p \geq 0, q > 0, s = p + q$, 求级数 $\sum_{n=1}^{\infty} \frac{1}{(sn - p)(sn + q)}$ 的和.\par
% p175-b19; source PDF page 175
\blockanchor{p175-b19}
\question{4}令 $f(x, y) = \begin{cases} \frac{1}{1 + (1 + xy)^{\frac{1}{y}}}, & 0 \leq x \leq 1, 0 < y \leq 1; \\ \frac{1}{1 + e^x}, & 0 \leq x \leq 1, y = 0. \end{cases}$ 求极限 $\lim_{y \to 0} \int_0^1 f(x, y) \, \mathrm{d}x$.\par
% p175-b20; source PDF page 175
\blockanchor{p175-b20}
\question{5}应用曲面面积公式, 计算半径为 R 的球体表面积.\par
% p175-b21; source PDF page 175
\blockanchor{p175-b21}
\question{6}应用积分号下积求积分, 计算 $I = \int_{0}^{1} \frac{x^{b} - x^{a}}{\ln x} \mathrm{d}x, b > a > 0$.\par
% p175-b22; source PDF page 175
\par\Needspace{5\baselineskip}
\blockanchor{p175-b22}
\majorquestion{二}简答题. 每题 10 分, 共 40 分.\par
% p175-b23; source PDF page 175
\blockanchor{p175-b23}
\question{1}求曲线  $\left\{\begin{aligned}x &= a(\cos t + t \sin t); \\ y &= a(\sin t - t \cos t).\end{aligned}\right.$  上任一点的法线到原点的距离.\par
% p175-b24; source PDF page 175
\blockanchor{p175-b24}
\question{2}设 $f_{n}(x) = \frac{x}{1 + n^{2}x^{2}}$，讨论函数序列 $\{f_{n}(x)\}$ 和 $\{f_{n}'(x)\}$ 在 $(- \infty, +\infty)$ 上的一致收敛性.\par
% p176-h00; source PDF page 176
\par\Needspace{6\baselineskip}
\blockanchor{p176-h00}
\question{3}计算\par
\nopagebreak[4]
% p176-b00; source PDF page 176
\blockanchor{p176-b00}
\displayformula{\iint_ {\Sigma} x \mathrm{d} y \mathrm{d} z + 2 y \mathrm{d} z \mathrm{d} x + z ^ {2} \mathrm{d} x \mathrm{d} y.}
% p176-b01; source PDF page 176
\blockanchor{p176-b01}
其中 $\Sigma$ 为椭球面 $\frac{x^2}{a^2} +\frac{y^2}{b^2} +\frac{z^2}{c^2} = 1$ 的内侧， $a,b,c$ 均为正数\par
% p176-b02; source PDF page 176
\par\Needspace{6\baselineskip}
\blockanchor{p176-b02}
\question{4}计算\par
\nopagebreak[4]
% p176-b03; source PDF page 176
\blockanchor{p176-b03}
\displayformula{I = \oint_ {C} \frac {e ^ {y}}{x ^ {2} + y ^ {2}} [ (x \sin x + y \cos x) \mathrm{d} x + (y \sin x - x \cos x) \mathrm{d} y ].}
% p176-b04; source PDF page 176
\blockanchor{p176-b04}
其中 $C: x^{2} + y^{2} = 1$, 沿逆时针方向.\par
% p176-b05; source PDF page 176
\par\Needspace{5\baselineskip}
\blockanchor{p176-b05}
\majorquestion{三}证明题. 每题10分, 共50分.\par
% p176-b06; source PDF page 176
\blockanchor{p176-b06}
\question{1}设 $\lim_{n\to \infty}a_n = a$ ，证明：\par
% p176-b07; source PDF page 176
\blockanchor{p176-b07}
(1) $\lim_{n\to \infty}\frac{[na_n]}{n} = a.$\par
% p176-b08; source PDF page 176
\blockanchor{p176-b08}
(2) 若 $a > 0, a_n > 0$, 则 $\lim_{n \to \infty} \sqrt[n]{a_n} = 1$.\par
% p176-b09; source PDF page 176
\blockanchor{p176-b09}
\question{2}证明不等式 $\frac{\tan x}{x} > \frac{x}{\sin x}, x \in \left(0, \frac{\pi}{2}\right)$.\par
% p176-b10; source PDF page 176
\blockanchor{p176-b10}
\question{3}设 $f(x)$ 为定义在 $[a, +\infty)$ 上的增函数. 证明: $\lim_{x \to +\infty} f(x)$ 存在的充要条件是 $f(x)$ 在 $[a, +\infty)$ 上有上界.\par
% p176-b11; source PDF page 176
\blockanchor{p176-b11}
\question{4}证明: 利用变量代换 $u = x^{2} - y^{2}$, $v = 2xy$ 可将方程 $\frac{\partial^2 z}{\partial x^2} + \frac{\partial^2 z}{\partial y^2} = 0$ 化为 $\frac{\partial^2 u}{\partial u^2} + \frac{\partial^2 z}{\partial v^2} = 0$.\par
% p176-b12; source PDF page 176
\blockanchor{p176-b12}
\question{5}若二元函数 $f(x, y)$ 在点 (0, 0) 处连续.\par
% p176-b13; source PDF page 176
\blockanchor{p176-b13}
(1) 证明: 若极限 $\lim_{(x,y)\to (0,0)}\frac{f(x,y)}{x^2 + y^2}$ 存在, 则函数 $f(x,y)$ 在点 $(0,0)$ 处可微.\par
% p176-b14; source PDF page 176
\blockanchor{p176-b14}
(2) 举例说明, 反之不成立.\par
% p176-b15; source PDF page 176
\par\Needspace{5\baselineskip}
\blockanchor{p176-b15}
\majorquestion{四}证明题. 共 15 分.\par
% p176-b16; source PDF page 176
\par\Needspace{6\baselineskip}
\blockanchor{p176-b16}
(1) 若 $f$ 和 $g$ 在 $[a, b]$ 上可积, 则\par
\nopagebreak[4]
% p176-b17; source PDF page 176
\blockanchor{p176-b17}
\displayformula{\left(\int_ {a} ^ {b} f (x) g (x) \mathrm{d} x\right) ^ {2} \leq \int_ {a} ^ {b} f ^ {2} (x) \mathrm{d} x \int_ {a} ^ {b} g ^ {2} (x) \mathrm{d} x.}
% p176-b18; source PDF page 176
\par\Needspace{6\baselineskip}
\blockanchor{p176-b18}
(2) 若 $f$ 在 $[a, b]$ 上可积, 则\par
\nopagebreak[4]
% p176-b19; source PDF page 176
\blockanchor{p176-b19}
\displayformula{\left(\int_ {a} ^ {b} f (x) \mathrm{d} x\right) ^ {2} \leq (b - a) \int_ {a} ^ {b} f ^ {2} (x) \mathrm{d} x.}
% p176-b20; source PDF page 176
\par\Needspace{5\baselineskip}
\blockanchor{p176-b20}
\majorquestion{五}证明题. 共 15 分. 设函数 $f(x)$ 在 $[0,1]$ 上连续, 在 $(0,1)$ 内可微, 且 $f(0) = 0$, $f(1) = 1$. 证明:\par
% p176-b21; source PDF page 176
\blockanchor{p176-b21}
(1) 存在 $\xi \in (0,1)$, 使得 $f(\xi) = 1 - \xi$.\par
% p176-b22; source PDF page 176
\blockanchor{p176-b22}
(2) 存在不同的两点 $\eta, \xi \in (0,1)$, 使得 $f'(\eta)f'(\xi) = 1$.\par
